% Part: normal-modal-logic
% Chapter: syntax-and-semantics
% Section: language-modal-logic

\documentclass[../../../include/open-logic-section]{subfiles}

\begin{document}

\olfileid[zh]{nml}{syn}{lan}

\olsection{基本模态逻辑的语言}

\begin{defn}
模态逻辑的基本语言包含
\begin{enumerate}
  \tagitem{prvFalse}{表示!!{falsity}的命题常元~$\lfalse$。}{}
  \tagitem{prvTrue}{表示!!{truth}的命题常元~$\ltrue$。}{}
  \item 一个!!{denumerable}的!!{propositional variable}集合：$\Obj
    p_0$, $\Obj p_1$, $\Obj p_2$, \dots
  \item 命题联结词：\startycommalist
  \iftag{prvNot}{\ycomma $\lnot$（否定）}{}%
  \iftag{prvAnd}{\ycomma $\land$（合取）}{}%
  \iftag{prvOr}{\ycomma $\lor$（析取）}{}%
  \iftag{prvIf}{\ycomma $\lif$（!!{conditional}）}{}%
  \iftag{prvIff}{\ycomma $\liff$（!!{biconditional}）}{}。
  \tagitem{prvBox}{模态算子 $\Box$。}{}
  \tagitem{prvDiamond}{模态算子 $\Diamond$。}{}
\end{enumerate}
\end{defn}

\begin{defn}
基本模态语言的\emph{!!^{formula}} 归纳定义如下：
\begin{enumerate}
\tagitem{prvFalse}{$\lfalse$ 是一个原子!!{formula}。}{}

\tagitem{prvTrue}{$\ltrue$ 是一个原子!!{formula}。}{}

\item 每个!!{propositional variable} $\Obj p_i$ 都是（原子）!!{formula}。

\tagitem{prvNot}{若 $!A$ 是!!a{formula}，则 $\lnot !A$ 是!!a{formula}。}{}

\tagitem{prvAnd}{若 $!A$ 和 $!B$ 是!!{formula}，则 $(!A \land
  !B)$ 是!!a{formula}。}{}

\tagitem{prvOr}{若 $!A$ 和 $!B$ 是!!{formula}，则 $(!A \lor !B)$
  是!!a{formula}。}{}

\tagitem{prvIf}{若 $!A$ 和 $!B$ 是!!{formula}，则 $(!A \lif !B)$
  是!!a{formula}。}{}

\tagitem{prvIff}{若 $!A$ 和 $!B$ 是!!{formula}，则 $(!A \liff !B)$
  是!!a{formula}。}{}

\tagitem{prvBox}{若 $!A$ 是!!a{formula}，则 $\Box !A$ 是!!a{formula}。}{}

\tagitem{prvDiamond}{若 $!A$ 是!!a{formula}，则 $\Diamond !A$ 是!!a{formula}。}{}

\tagitem{limitClause}{除此之外，没有其他!!a{formula}。}{}
\end{enumerate}
\end{defn}

\begin{tagblock}{defNot,defOr,defAnd,defIf,defIff,defTrue,defFalse,defBox,defDiamond}
\begin{defn}
用已定义的算子构造的公式按如下方式理解：

\begin{tagenumerate}{defTrue,defFalse,defNot,defOr,defAnd,defIf,defIff,defBox,defDiamond}

\tagitem{defTrue}{$\ltrue$ 缩写为
  \iftag{prvFalse}{$\lnot\lfalse$}{$(!A \lor \lnot !A)$，其中 $!A$ 为
    某个固定的原子!!{formula}~$!A$}。}{}

\tagitem{defFalse}{$\lfalse$ 缩写为
  \iftag{prvTrue}{$\lnot\ltrue$}{$(!A \land \lnot !A)$，其中 $!A$ 为
    某个固定的原子!!{formula}~$!A$}。}{}

\tagitem{defNot}{$\lnot !A$ 缩写为 $!A \lif \lfalse$。}{}

\tagitem{defOr}{$!A \lor !B$ 缩写为
  \iftag{prvAnd}{$\lnot(\lnot !A \land \lnot !B)$}{$\lnot !A \lif
    !B$}。}{}

\tagitem{defAnd}{$!A \land !B$ 缩写为
  \iftag{prvOr}{$\lnot(\lnot !A \lor \lnot !B)$}{$\lnot (!A \lif
    \lnot !B)$}。}{}

\tagitem{defIf}{$!A \lif !B$ 缩写为
  \iftag{prvOr}{$\lnot !A \lor !B)$}{$\lnot (!A \land \lnot !B)$}。}{}

\tagitem{defIff}{$!A \liff !B$ 缩写为 $(!A \lif !B) \land (!B
  \lif !A)$。}{}

\tagitem{defBox}{$\Box !A$ 缩写为 $\lnot\Diamond\lnot !A$。}{}

\tagitem{defDiamond}{$\Diamond !A$ 缩写为 $\lnot\Box\lnot !A$。}{}
\end{tagenumerate}
\end{defn}
\end{tagblock}

若!!a{formula}~$!A$ 不含 $\Box$ 或 $\Diamond$，我们称它为
\emph{无模态的}。

\end{document}
